% GENERATED FILE. Edit canonical lessons, then run npm run generate:books.
% canonical-source-hash: fnv1a64:4fbd2d082d075233
% canonical-lessons: LA-C548-pulchritudo, LA-C548-speciosus, LA-C548-elegans, LA-C548-ornamentum, LA-C548-deformis

\chapter{Beauty}
\label{ch:la-beauty}
\hlchaptermodality{548}

\begin{chapteropening}
\textbf{By the end of this chapter:} \emph{I can describe beauty, charm and ugliness.}

Complete the last lesson of chapter 548: I can describe beauty, charm and ugliness.
\end{chapteropening}

\section[pulchritūdō, pulchritūdinis]{pulchritūdō, pulchritūdinis — beauty}

\label{lesson:LA-C548-pulchritudo}



\begin{quote}
Before the new one: say the Latin for a small cube, then the Latin for to engrave.
\end{quote}

\subsection*{You'll want to know: pulchritūdō, pulchritūdinis}
\textbf{pulchritūdō, pulchritūdinis} — \textquotedblleft{}beauty\textquotedblright{}.

\begin{quote}
Pulchritūdō hortī hospitēs dēlectat.
\end{quote}

\begin{quote}
\emph{The beauty of the garden delights the guests.}
\end{quote}

Feminine, third declension, from \emph{pulcher}. English \emph{pulchritude}.

\begin{cousinweb}
\textbf{pulchritūdō} is built on \emph{pulcher} \textquotedblleft{}beautiful\textquotedblright{} with \emph{-tūdō} (EtymDB; Lewis \& Short). English \emph{pulchritude} comes from it.
\end{cousinweb}

\begin{grammarlens}[title={What you've built}]
One more word for telling, explaining and describing at B1.
\end{grammarlens}

\subsection*{Your turn}
\begin{itemize}
\raggedright
  \item \emph{Say it:} \emph{pulchritūdō, pulchritūdinis}
  \item \emph{Say it:} \emph{pulchritūdō, pulchritūdinis}, once more
  \item \emph{Say it:} a sentence of your own with \emph{pulchritūdō, pulchritūdinis}
  \item \emph{Recall:} say the Latin for a small cube, then the Latin for to engrave, then say \emph{pulchritūdō, pulchritūdinis} again
\end{itemize}

\begin{tcolorbox}[breakable,colback=teal!4,colframe=teal!35!black,title={Before you move on}]
What does pulchritūdō, pulchritūdinis mean? (\textquotedblleft{}Beauty\textquotedblright{}.) Say it once more.
\end{tcolorbox}
\section[speciōsus, speciōsa, speciōsum]{speciōsus, speciōsa, speciōsum — good-looking, handsome}

\label{lesson:LA-C548-speciosus}



\begin{quote}
Before the new one: say the Latin for to engrave, then the Latin for beauty.
\end{quote}

\subsection*{You'll want to know: speciōsus, speciōsa, speciōsum}
\textbf{speciōsus, speciōsa, speciōsum} — \textquotedblleft{}good-looking, handsome\textquotedblright{}.

\begin{quote}
Iuvenis speciōsus est, sed superbus.
\end{quote}

\begin{quote}
\emph{The young man is good-looking, but proud.}
\end{quote}

From \emph{speciēs} — appearance. Later also showy but false: English \emph{specious}.

\begin{cousinweb}
\textbf{speciōsus} is built on \emph{speciēs} \textquotedblleft{}appearance, look\textquotedblright{}, from \emph{speciō} \textquotedblleft{}to behold\textquotedblright{} (Lewis \& Short). English \emph{specious} comes from it, with the later sense \textquotedblleft{}fair on the surface\textquotedblright{}.
\end{cousinweb}

\begin{grammarlens}[title={What you've built}]
One more word for telling, explaining and describing at B1.
\end{grammarlens}

\subsection*{Your turn}
\begin{itemize}
\raggedright
  \item \emph{Say it:} \emph{speciōsus, speciōsa, speciōsum}
  \item \emph{Say it:} \emph{speciōsus, speciōsa, speciōsum}, once more
  \item \emph{Say it:} a sentence of your own with \emph{speciōsus, speciōsa, speciōsum}
  \item \emph{Recall:} say the Latin for to engrave, then the Latin for beauty, then say \emph{speciōsus, speciōsa, speciōsum} again
\end{itemize}

\begin{tcolorbox}[breakable,colback=teal!4,colframe=teal!35!black,title={Before you move on}]
What does speciōsus, speciōsa, speciōsum mean? (\textquotedblleft{}Good-looking, handsome\textquotedblright{}.) Say it once more.
\end{tcolorbox}
\section[ēlegāns, ēlegantis]{ēlegāns, ēlegantis — elegant, refined, tasteful}

\label{lesson:LA-C548-elegans}



\begin{quote}
Before the new one: say the Latin for beauty, then the Latin for good-looking.
\end{quote}

\subsection*{You'll want to know: ēlegāns, ēlegantis}
\textbf{ēlegāns, ēlegantis} — \textquotedblleft{}elegant, refined, tasteful\textquotedblright{}.

\begin{quote}
Domus eius parva, sed ēlegāns est.
\end{quote}

\begin{quote}
\emph{His house is small but elegant.}
\end{quote}

One ending. First choosy, then of good taste. English \emph{elegant}.

\begin{cousinweb}
\textbf{ēlegāns} is probably an old form of \emph{ēligēns}, the participle of \emph{ēligō} \textquotedblleft{}to pick out\textquotedblright{} (Lewis \& Short; EtymDB), from \emph{legō} \textquotedblleft{}to choose\textquotedblright{}. English \emph{elegant} comes from it.
\end{cousinweb}

\begin{grammarlens}[title={What you've built}]
One more word for telling, explaining and describing at B1.
\end{grammarlens}

\subsection*{Your turn}
\begin{itemize}
\raggedright
  \item \emph{Say it:} \emph{ēlegāns, ēlegantis}
  \item \emph{Say it:} \emph{ēlegāns, ēlegantis}, once more
  \item \emph{Say it:} a sentence of your own with \emph{ēlegāns, ēlegantis}
  \item \emph{Recall:} say the Latin for beauty, then the Latin for good-looking, then say \emph{ēlegāns, ēlegantis} again
\end{itemize}

\begin{tcolorbox}[breakable,colback=teal!4,colframe=teal!35!black,title={Before you move on}]
What does ēlegāns, ēlegantis mean? (\textquotedblleft{}Elegant, refined, tasteful\textquotedblright{}.) Say it once more.
\end{tcolorbox}
\section[ōrnāmentum, ōrnāmentī]{ōrnāmentum, ōrnāmentī — an ornament, decoration}

\label{lesson:LA-C548-ornamentum}



\begin{quote}
Before the new one: say the Latin for good-looking, then the Latin for elegant.
\end{quote}

\subsection*{You'll want to know: ōrnāmentum, ōrnāmentī}
\textbf{ōrnāmentum, ōrnāmentī} — \textquotedblleft{}an ornament, decoration\textquotedblright{}.

\begin{quote}
Domina ōrnāmenta aurea gerit.
\end{quote}

\begin{quote}
\emph{The lady wears gold ornaments.}
\end{quote}

Neuter, second declension, from \emph{ōrnō} — to equip, adorn. English \emph{ornament}.

\begin{cousinweb}
\textbf{ōrnāmentum} is built on \emph{ōrnō} \textquotedblleft{}to equip, adorn\textquotedblright{} with \emph{-mentum} (EtymDB; Lewis \& Short). English \emph{ornament} comes from it.
\end{cousinweb}

\begin{grammarlens}[title={What you've built}]
One more word for telling, explaining and describing at B1.
\end{grammarlens}

\subsection*{Your turn}
\begin{itemize}
\raggedright
  \item \emph{Say it:} \emph{ōrnāmentum, ōrnāmentī}
  \item \emph{Say it:} \emph{ōrnāmentum, ōrnāmentī}, once more
  \item \emph{Say it:} a sentence of your own with \emph{ōrnāmentum, ōrnāmentī}
  \item \emph{Recall:} say the Latin for good-looking, then the Latin for elegant, then say \emph{ōrnāmentum, ōrnāmentī} again
\end{itemize}

\begin{tcolorbox}[breakable,colback=teal!4,colframe=teal!35!black,title={Before you move on}]
What does ōrnāmentum, ōrnāmentī mean? (\textquotedblleft{}An ornament, decoration\textquotedblright{}.) Say it once more.
\end{tcolorbox}
\section[dēfōrmis, dēfōrme]{dēfōrmis, dēfōrme — ugly, misshapen}

\label{lesson:LA-C548-deformis}



\begin{quote}
Before the new one: say the Latin for elegant, then the Latin for an ornament.
\end{quote}

\subsection*{You'll want to know: dēfōrmis, dēfōrme}
\textbf{dēfōrmis, dēfōrme} — \textquotedblleft{}ugly, misshapen\textquotedblright{}.

\begin{quote}
Statua dēfōrmis est, sed antīqua.
\end{quote}

\begin{quote}
\emph{The statue is ugly, but old.}
\end{quote}

Two endings. \emph{dē-} 'away from' + \emph{fōrma} 'shape, beauty'.

\begin{cousinweb}
\textbf{dēfōrmis} is built on \emph{dē-} \textquotedblleft{}away from\textquotedblright{} and \emph{fōrma} \textquotedblleft{}shape, beauty\textquotedblright{} (Lewis \& Short; EtymDB). English \emph{deformed} comes from the related verb \emph{dēfōrmō}.
\end{cousinweb}

\begin{grammarlens}[title={What you've built}]
One more word for telling, explaining and describing at B1.
\end{grammarlens}

\subsection*{Your turn}
\begin{itemize}
\raggedright
  \item \emph{Say it:} \emph{dēfōrmis, dēfōrme}
  \item \emph{Say it:} \emph{dēfōrmis, dēfōrme}, once more
  \item \emph{Say it:} a sentence of your own with \emph{dēfōrmis, dēfōrme}
  \item \emph{Recall:} say the Latin for elegant, then the Latin for an ornament, then say \emph{dēfōrmis, dēfōrme} again
\end{itemize}

\begin{tcolorbox}[breakable,colback=teal!4,colframe=teal!35!black,title={Before you move on}]
What does dēfōrmis, dēfōrme mean? (\textquotedblleft{}Ugly, misshapen\textquotedblright{}.) Say it once more.
\end{tcolorbox}
